\documentclass[../main.tex]{subfiles}

\begin{document}

\chapter{不变测度}\label{chapter:ch22-不变测度}

一个拓扑群是群 $G$ 与 $G$ 上的 Hausdorff 拓扑一起使得群运算和逆是连续的. 我们证明 John von Neumann 的一个重要定理, 它告诉我们在任何紧拓扑群 $G$ 上存在唯一的测度 $\mu$, 称为 Haar 测度, 它在群的作用下不变, 即
\begin{align*}
\text{对所有 }g\in G,\ E\in\mathcal{B}(G),\quad\mu(g\cdot E)=\mu(E).
\end{align*}
唯一性从 Fubini 定理得出; 存在性是 Shizuo Kakutani 不动点定理的一个推论, 它断言对一个紧群 $G$, 存在泛函 $\psi\in[C(G)]^{*}$ 使得对所有 $g\in G$, $f\in C(G)$,
\begin{align*}
\psi[f\equiv 1]=1\text{ 且 }\psi[x\mapsto f(x)]=\psi[x\mapsto f(g\cdot x)].
\end{align*}
在这个不动点定理的证明中 Alaoglu 定理是关键的. Haar 测度的存在性的证明细节在 $G$ 到 $(C(G))^{*}$ 的一般线性群的群同态背景下构建. 我们也考虑紧度量空间 $X$ 到其自身的映射 $f$ 与 $\mathcal{B}(X)$ 上的有限测度. 基于 Helly 定理, 我们证明了 Bogoliubov-Krilov 定理. 它告诉我们, 若 $f$ 是紧度量空间 $X$ 上的连续映射, 则存在 $\mathcal{B}(X)$ 上的测度 $\mu$ 满足
\begin{align*}
\mu(X)=1\text{ 且对所有 }\varphi\in L^{1}(X,\mu),\quad\int_{X}\varphi\circ f\,d\mu=\int_{X}\varphi\,d\mu.
\end{align*}
基于 Krein-Milman 定理, 我们证明了上述的 $\mu$ 可以恰当选取使得 $f$ 关于 $\mu$ 是遍历的, 即若 $A$ 属于 $\mathcal{B}(X)$ 且 $\mu([A\sim f(A)]\cup[f(A)\sim A])=0$, 则 $\mu(A)=0$ 或 $\mu(A)=1$.

\subfile{Chapter-22/section-01.tex}
\subfile{Chapter-22/section-02.tex}
\subfile{Chapter-22/section-03.tex}
\subfile{Chapter-22/section-04.tex}

\end{document}